\section{Supporting Interaction-Field Results and Implementation}
\label{sec:supporting-results}

This appendix documents the sample roster, the resampling and inference
architecture, the transport-dimension ladder, the representation sensitivity,
and the annual joint-field estimates.

The firms used to construct the primary barycentric interaction field are
listed in \Cref{tab:sample-universe}.

% AUTO-GENERATED by pipeline render-sample-universe-table -- do not edit by hand.

\begin{table}[htbp]
\centering
\small
\begin{tabularx}{\linewidth}{@{}lr>{\raggedright\arraybackslash}X@{}}
\toprule
Sector & Firms & Ticker symbols \\
\midrule
Communication Services & 11 & BIDU, CHTR, CMCSA, EA, GOOG, MTCH, NFLX, NTES, SIRI, TMUS, TTWO \\
Consumer Cyclical & 18 & BKNG, EBAY, EXPE, HAS, JD, LULU, MAR, MAT, MELI, NCLH, ORLY, ROST, SBUX, TCOM, TRIP, TSCO, ULTA, WYNN \\
Consumer Defensive & 7 & COST, DLTR, KDP, KHC, MDLZ, MNST, PEP \\
Energy & 1 & FANG \\
Financial Services & 1 & PYPL \\
Healthcare & 14 & ALGN, AMGN, AZN, BIIB, BMRN, DXCM, GILD, HOLX, IDXX, ILMN, INCY, ISRG, REGN, VRTX \\
Industrials & 7 & AAL, CSX, FAST, HON, JBHT, PCAR, UAL \\
Technology & 38 & ADBE, ADI, ADP, ADSK, AKAM, AMAT, AMD, ASML, AVGO, CDNS, CHKP, CSCO, CTSH, ENPH, FISV, FTNT, INTC, INTU, KLAC, LRCX, MCHP, MRVL, MU, NTAP, NVDA, NXPI, OKTA, PANW, PAYX, QCOM, SNPS, STX, SWKS, TEAM, TXN, WDAY, WDC, ZS \\
Utilities & 3 & AEP, EXC, XEL \\
\bottomrule
\end{tabularx}
\caption{Canonical empirical roster: ticker symbols of the firms in the canonical priced $W_2$ geometry, grouped by the sector label in Nasdaq summary metadata. Firm names accompany the symbols in the replication archive's universe metadata.}
\label{tab:sample-universe}
\end{table}


\subsection{Resampling and Inference Implementation}
\label{sec:app-inference}\label{sec:p5-two-field-implementation}

Every interval and every held-out comparison in the paper draws on one
resampling architecture, so a gain reported in more than one table carries one
interval.
Each of the 2,000 draws uses three stationary-block date schedules with
expected block length 21, taken from a single seed: one for the 2018--2022
calibration panel on which the scored parameters are fitted, one for the
post-corpus panel on which in-period coefficients are refitted, and one for the
post-corpus dates over which the fitted model is scored.
Every model, every comparator field, and both transport orders share the same
draws, so score differences are paired, and each draw preserves the
cross-section observed on its resampled dates.
The procedure refits the factor-only restriction, the direct Ledoit--Wolf
residual covariance, and every field coefficient on each calibration draw; the
calibration-tuned RBF bandwidth is held at its full-calibration selection, so
that operator is frozen while its coefficient is refitted.
Annual folds add one to the base seed in period order, and within a period the
joint fit and both single-field boundaries share one bootstrap bank.
The draws condition on the corpus, embeddings, and constructed fields, so they
propagate return-sampling and nuisance-parameter uncertainty but not
uncertainty from article sampling, from the representation, or from tied
optimal assignments.
All reported intervals are bootstrap quantiles; numerical-Hessian standard
errors are retained only as diagnostics, and boundary-constrained joint
intervals are descriptive.

The joint search profiles the mixture weight and total feedback on a
\ppvalue{two-field-theta-points} by \ppvalue{two-field-rho-points} grid and
refines the maximum continuously within its grid cell; every bootstrap draw
receives the same refinement, so coefficients and intervals do not inherit the
grid spacing.
The Jacobian term is evaluated from the eigenvalues of each field mixture on
the grid and by direct LU factorization during refinement; tests verify that
the two routes agree to floating-point tolerance.

The restricted-null QLR calibration is used only for the nested restrictions
$\rho=0$, $\rho^\flat=0$, and $\rho^{\mathrm{news}}=0$, never for the
non-nested comparison
between fields.
For each null, the procedure recovers and centers the residuals of the
restricted fit, standardizes each asset path, and draws an independent
stationary-block path for every asset and null replicate, which preserves
temporal blocks while removing contemporaneous residual alignment under the
working null.
After restoring the fitted common or firm-specific variances and factor means,
it maps the innovations back through the restricted spatial inverse and refits
both models with the same refinement; non-converged draws are rejected.
The $p$-value is $(1+\#\{QLR^*\geq QLR\})/(B+1)$ with $B=2{,}000$ null draws,
reported at its floor of $1/2{,}001$ when no draw exceeds the observed
statistic.
The simulation accommodates the nonnegative-channel boundary, so no chi-square
reference is used; this working null does not provide size control under
common stochastic volatility.

The transport-order equivalence margin, fixed before the final estimation, is
1\% of the W2-versus-factor-only mean calibration-score gain: the fields are
predictively equivalent when the 90\% paired interval for the W2-minus-W1
score lies inside it.
The weak-separation rule, fixed after a pilot computation and before the same
final estimation, classifies the fields as nearly collinear when their induced
peer-return signals correlate at $0.98$ or more and the standard deviation of
their difference is at most $0.10$ of the W2 signal's.

\subsection{Matryoshka Dimension Ladder}
\label{sec:p5-matryoshka}

The ladder repeats both transport orders at $D\in\{4096,1024,256,64\}$ within
Qwen3-Embedding 8B \citep{kusupati_matryoshka_2022,qwen3_embedding_2025}.
For each article embedding $z$, the representation retains the leading $D$
coordinates and then normalizes that prefix,
\[
  z^{(D)}=\frac{(z_1,\ldots,z_D)}{\lVert(z_1,\ldots,z_D)\rVert_2}.
\]
The ground cost remains unscaled Euclidean chord distance at every dimension.
Unit normalization bounds chord distance in $[0,2]$ and squared chord cost in
$[0,4]$, and for two unit prefixes $d^2=2-2\langle z_i^{(D)},z_j^{(D)}\rangle$,
so the W2 assignment cost is twice their cosine dissimilarity.
A common scalar rescaling multiplies the W1 and W2 assignment objectives by
positive constants, so it alters neither optimal assignment nor the minimizing
reconstruction weights.

The mechanism under study is concentration rather than scale.
For any reference distance $d_0$, the exact identity
\[
  d^2=2d_0d-d_0^2+(d-d_0)^2
\]
shows that squared cost is affine in cost up to the final nonnegative
remainder.
When the observed distances occupy a narrow neighborhood of $d_0$, that
remainder is small and minimum-sum $d$ and $d^2$ assignments can coincide.
The ladder tests this local approximation empirically.
At each dimension it enumerates every article-to-article cost in every
unordered cross-firm balanced assignment problem, reports the population
coefficient of variation of chord distance, and traces W1 and W2 agreement
through the share of identical atom assignments, the correlation of
off-diagonal operator weights, the correlation of induced peer-return signals,
and the W2-minus-W1 held-out score under the primary specification.
This ordering separates concentration in the embedding geometry from
convergence introduced by target reconstruction or by common variation in
returns.

% AUTO-GENERATED -- do not edit by hand.
\begin{table}[H]
\centering
\scriptsize
\setlength{\tabcolsep}{3.2pt}
\begin{tabular}{rrrrrr}
\toprule
Dim. & Chord-distance CV & Assignment overlap & Weight corr. & Spatial-lag corr. & $\Delta$ held-out score \\
\midrule
4,096 & 5.75\% & 90.1\% & 0.9979 & 0.9998 & -0.000109 \\
1,024 & 5.97\% & 89.5\% & 0.9978 & 0.9998 & -0.000339 \\
256 & 7.16\% & 87.4\% & 0.9977 & 0.9998 & -0.000354 \\
64 & 10.33\% & 82.4\% & 0.9975 & 0.9997 & -0.000443 \\
\bottomrule
\end{tabular}
\caption{Matryoshka dimension ladder for target-anchored W1 and W2 transport. At each dimension, leading coordinates are truncated and then row-wise L2-normalized; the ground cost is unscaled Euclidean chord distance. Chord-distance CV uses all 81,100,800 article-pair costs across 4,950 unordered firm pairs. Assignment overlap is the share of the 633,600 matched target atoms receiving the same peer atom under minimum-sum $d$ and $d^2$ assignments. Weight correlation uses off-diagonal target-operator entries; spatial-lag correlation uses the common 2023--2026 excess-return panel. The final column is W2 minus W1 mean Gaussian quasi-log score per firm-day after fitting the FF5-plus-momentum model on 2018--2022; positive values favour W2.}
\label{tab:transport-dimension-ladder}
\end{table}


\subsection{Embedding-Representation Sensitivity}
\label{sec:p5-representation-ablation}

The representation exercises ask whether the field comparison depends on one
text representation.
The sensitivity roster separates four design margins: output width, model
capacity, model family, and information vintage.
An encoder maps each article into a vector, whose dimension determines the
coordinates entering both the pairwise $W_2$ distances and the target-anchored
reconstruction.

The output-width comparison uses Qwen3-Embedding 8B at full width and at 1,024,
256, and 64 coordinates.
Qwen's Matryoshka training is designed to keep leading coordinate prefixes
informative, so the shorter vectors test how much compression the interaction
geometry tolerates.
The common 1,024-coordinate Qwen3-Embedding 8B and 4B rows then vary model
capacity while holding output width fixed.
Their full-width comparison changes capacity and output width together, so it
does not isolate capacity.
The 1,024-coordinate BGE-large-v1.5 row changes model family and pretraining
relative to Qwen3-Embedding 8B, making it an external sensitivity check rather
than an identified training-mechanism effect
\citep{qwen3_embedding_2025,xiao_cpack_2023}.

The primary Qwen encoder postdates the 2018--2022 article window, so the
information-vintage comparison asks whether the field comparison depends on an
encoder's information cutoff.
It uses the published ChronoBERT checkpoint for each year-end from 2017 through
2024 \citep{he_chronobert_2025}.
These ModernBERT-based masked-language encoders produce 768-coordinate
representations by averaging final-layer token states over all non-padding
positions, including special tokens.
We hold inference fixed at 8,000 characters and 1,024 tokens per article,
using bfloat16 autocast with float32 model parameters and token-mean pooling.
We apply row-L2 normalization once before constructing distances.
The same 128 articles per firm enter every checkpoint; neither layers nor
checkpoints are selected using returns.
Later checkpoints continue training their predecessors, so their contrasts
combine later information with cumulative training exposure, not matched
compute or an identified causal leakage effect.
The 2017 cutoff precedes article construction; 2018--2022 falls within it,
and 2023--2024 postdates construction but overlaps the full 2023--2026
return-evaluation period.
No arm is a post-evaluation negative control.
Since publication occurred later, the exercise concerns reported training
cutoffs, not historical deployment availability or point-in-time validity
of all other study inputs.

Every row in \Cref{tab:p5-representation-ablation} uses the same article inputs
through 2022, the same 2023--2026 return dates, and one common joint-date
bootstrap schedule.
Pairing each contrast within the same resample removes common bootstrap
variation, while the Holm correction accounts for testing several
representations in each family.
The table reports direct differences in $\rho$ for the barycentric field and
changes in its $\rho$ gap relative to RBF.
For representation $r$, write
$g_r=\rho^\flat_r-\rho^h_r$.
The two candidate-minus-reference estimands are
$\rho^\flat_c-\rho^\flat_r$ and $g_c-g_r$; negative values indicate,
respectively, weaker fitted feedback for the barycentric field and a smaller
advantage over the matching RBF operator.
The persistent co-mention graph is embedding-free and is therefore reported
once.
The equal-active-support comparator replaces the fitted coordinate magnitudes
with uniform weights on the same selected support, so it asks whether the
barycentric coordinate magnitudes improve conditional spatial fit beyond peer
selection.

\subsection{Point-in-Time Encoder Check}
\label{sec:p5-point-in-time}

The primary Qwen encoder was released after the article-construction window.
This sensitivity therefore repeats the primary FF5-plus-momentum held-out
score comparison with ChronoBERT 2017, whose reported training cutoff is
31 December 2017, before both article construction and calibration
\citep{he_chronobert_2025}.
The barycentric and median-bandwidth RBF fields are rebuilt from the pre-period
representation, held fixed during evaluation, and compared with the same
factor-only restriction.

% AUTO-GENERATED -- do not edit by hand.
\begin{table}[H]
\centering
\scriptsize
\setlength{\tabcolsep}{3.2pt}
\begin{tabular}{lrrrr}
\toprule
Interaction field & $\hat\rho$ & Held-out score & Gain & 95\% CI \\
\midrule
Factor-only & 0.000 & 2.468881 & 0.000000 & -- \\
ChronoBERT 2017 W2 barycentric & 0.561 & 2.497084 & 0.028203 & [0.0249, 0.0318] \\
ChronoBERT 2017 median-bandwidth W2 RBF & 0.452 & 2.478885 & 0.010004 & [0.0082, 0.0120] \\
\midrule
W2 minus median RBF & -- & -- & 0.018199 & [0.0152, 0.0214] \\
\bottomrule
\end{tabular}
\caption{Point-in-time representation sensitivity under the primary FF5-plus-momentum specification. ChronoBERT 2017 has a reported training-information cutoff of 2017-12-31, pre-period relative to the 2018--2022 article and calibration windows. This is not a claim of historical checkpoint availability. Fields are held fixed throughout the 2023--2026 evaluation. The score is the mean Gaussian quasi-log score per firm-day; intervals use paired stationary-block bootstrap draws. The RBF field uses the median off-diagonal squared W2 distance as its bandwidth.}
\label{tab:p5-point-in-time-encoder}
\end{table}


The pre-period representation preserves a positive barycentric gain
and a positive
paired barycentric-minus-RBF difference under the primary score design.

The contrast table states one joint equivalence rule for the vintage rows and
reports its two benchmark-calibrated bounds.
A pair meets the rule only when the marginal 95\% paired intervals for both
estimands lie strictly inside their corresponding bounds.
Holm adjusts the $p$-values separately for the 28 vintage comparisons in each
metric; the intervals are not simultaneous.
These post-specified thresholds organize a descriptive sensitivity exercise;
they are not prespecified confirmatory equivalence margins.

% AUTO-GENERATED -- do not edit by hand.
\begingroup
\scriptsize
\setlength{\tabcolsep}{3.3pt}
\begin{longtable}{@{}p{0.35\linewidth}llrrr@{}}
\caption{Pooled 2023--2026 representation sensitivity of the single-field spatial horse race across seven base representations and eight published annual ChronoBERT vintages. Each encoder row compares the barycentric interaction field $W^\flat$, obtained by target-anchored Wasserstein-barycentric reconstruction, with the median-bandwidth RBF transform of its matching squared-W$_2$ matrix. Intervals use the same 2,000 joint-date stationary bootstrap on the common 815-date schedule. $\hat\lambda=\hat\rho/(1-\hat\rho)$ is the implied adjustment index; log-likelihood is descriptive. Persistent news co-mentions are embedding-free; equal active support is shown once for the primary barycentric interaction field. These are paired conditional-fit sensitivity comparisons across the registered representation cells. Published cumulative-training checkpoints are not a matched-compute causal leakage test.}
\label{tab:p5-representation-ablation} \\
\toprule
Encoder & Dim. & Interaction matrix & $\hat\rho$ [95\% CI] & $\hat\lambda$ & Log-likelihood \\
\midrule
\endfirsthead
\multicolumn{6}{c}{\tablename\ \thetable\ (continued)} \\
\toprule
Encoder & Dim. & Interaction matrix & $\hat\rho$ [95\% CI] & $\hat\lambda$ & Log-likelihood \\
\midrule
\endhead
\midrule
\endfoot
\bottomrule
\endlastfoot
Qwen3-8B full (primary) & 4096 & Barycentric $W^\flat$ & 0.843 [0.815, 0.869] & 5.36 & 205146.3 \\
 &  & RBF--Wasserstein $W^h$ & 0.812 [0.780, 0.842] & 4.33 & 200099.7 \\
\addlinespace[1.5pt]
Qwen3-4B full & 2560 & Barycentric $W^\flat$ & 0.844 [0.816, 0.870] & 5.42 & 204835.4 \\
 &  & RBF--Wasserstein $W^h$ & 0.812 [0.779, 0.842] & 4.32 & 200068.5 \\
\addlinespace[1.5pt]
Qwen3-8B@1024 & 1024 & Barycentric $W^\flat$ & 0.844 [0.815, 0.870] & 5.40 & 205070.3 \\
 &  & RBF--Wasserstein $W^h$ & 0.812 [0.780, 0.842] & 4.33 & 200104.7 \\
\addlinespace[1.5pt]
Qwen3-4B@1024 & 1024 & Barycentric $W^\flat$ & 0.842 [0.813, 0.868] & 5.33 & 204859.1 \\
 &  & RBF--Wasserstein $W^h$ & 0.812 [0.780, 0.842] & 4.33 & 200094.3 \\
\addlinespace[1.5pt]
BGE-large-v1.5 full & 1024 & Barycentric $W^\flat$ & 0.840 [0.811, 0.866] & 5.24 & 204072.1 \\
 &  & RBF--Wasserstein $W^h$ & 0.811 [0.778, 0.841] & 4.29 & 200042.4 \\
\addlinespace[1.5pt]
Qwen3-8B@256 & 256 & Barycentric $W^\flat$ & 0.839 [0.810, 0.866] & 5.22 & 204870.9 \\
 &  & RBF--Wasserstein $W^h$ & 0.814 [0.781, 0.844] & 4.37 & 200127.5 \\
\addlinespace[1.5pt]
Qwen3-8B@64 & 64 & Barycentric $W^\flat$ & 0.821 [0.787, 0.852] & 4.59 & 204193.7 \\
 &  & RBF--Wasserstein $W^h$ & 0.816 [0.784, 0.845] & 4.43 & 200215.3 \\
\addlinespace[1.5pt]
ChronoBERT 2017 & 768 & Barycentric $W^\flat$ & 0.784 [0.741, 0.822] & 3.62 & 201973.8 \\
 &  & RBF--Wasserstein $W^h$ & 0.811 [0.777, 0.841] & 4.28 & 200041.7 \\
\addlinespace[1.5pt]
ChronoBERT 2018 & 768 & Barycentric $W^\flat$ & 0.781 [0.738, 0.820] & 3.56 & 201824.1 \\
 &  & RBF--Wasserstein $W^h$ & 0.810 [0.777, 0.841] & 4.28 & 200031.9 \\
\addlinespace[1.5pt]
ChronoBERT 2019 & 768 & Barycentric $W^\flat$ & 0.781 [0.738, 0.821] & 3.57 & 201786.3 \\
 &  & RBF--Wasserstein $W^h$ & 0.810 [0.776, 0.841] & 4.27 & 200026.7 \\
\addlinespace[1.5pt]
ChronoBERT 2020 & 768 & Barycentric $W^\flat$ & 0.781 [0.738, 0.820] & 3.57 & 201811.2 \\
 &  & RBF--Wasserstein $W^h$ & 0.810 [0.776, 0.841] & 4.27 & 200027.5 \\
\addlinespace[1.5pt]
ChronoBERT 2021 & 768 & Barycentric $W^\flat$ & 0.776 [0.732, 0.816] & 3.47 & 201664.2 \\
 &  & RBF--Wasserstein $W^h$ & 0.810 [0.776, 0.840] & 4.26 & 200027.8 \\
\addlinespace[1.5pt]
ChronoBERT 2022 & 768 & Barycentric $W^\flat$ & 0.782 [0.739, 0.820] & 3.58 & 201883.2 \\
 &  & RBF--Wasserstein $W^h$ & 0.810 [0.777, 0.841] & 4.27 & 200040.1 \\
\addlinespace[1.5pt]
ChronoBERT 2023 & 768 & Barycentric $W^\flat$ & 0.786 [0.744, 0.824] & 3.67 & 202016.8 \\
 &  & RBF--Wasserstein $W^h$ & 0.810 [0.777, 0.840] & 4.26 & 200043.2 \\
\addlinespace[1.5pt]
ChronoBERT 2024 & 768 & Barycentric $W^\flat$ & 0.783 [0.741, 0.822] & 3.62 & 201936.1 \\
 &  & RBF--Wasserstein $W^h$ & 0.810 [0.777, 0.840] & 4.27 & 200036.6 \\
\addlinespace[1.5pt]
\midrule
Embedding-free & -- & Persistent news $W^{\mathrm{news}}$ & 0.667 [0.618, 0.716] & 2.01 & 203493.8 \\
Primary support & -- & Equal active support & 0.841 [0.814, 0.865] & 5.27 & 202124.6 \\
\end{longtable}
\endgroup
\begingroup
\scriptsize
\setlength{\tabcolsep}{3.0pt}
\begin{longtable}{@{}p{0.35\linewidth}rrrrrr@{}}
\caption{Paired common-schedule QMLE contrasts. Panel A estimates $\rho^\flat_c-\rho^\flat_r$; Panel B estimates $(\rho^\flat_c-\rho^h_c)-(\rho^\flat_r-\rho^h_r)$. $p$ denotes two-sided add-one sign-tail probabilities; $p_H$ denotes Holm adjustment separately within each metric's seven base and 28 vintage contrasts. Estimates and intervals retain five decimal places so endpoints near zero remain visible. Vintage equivalence requires both intervals to lie strictly inside their respective BGE-large-minus-Qwen3-8B@1024 benchmark-calibrated, post-specified bounds; Eq. is this joint two-metric decision. Intervals are marginal, not simultaneous. The contrasts compare conditional fit on a common return schedule.}
\label{tab:p5-representation-contrasts} \\
\toprule
Contrast ($c-r$) & Estimate & 95\% CI & $p$ & $p_H$ & Bound & Eq. \\
\midrule
\endfirsthead
\multicolumn{7}{c}{\tablename\ \thetable\ (continued)} \\
\toprule
Contrast ($c-r$) & Estimate & 95\% CI & $p$ & $p_H$ & Bound & Eq. \\
\midrule
\endhead
\midrule
\endfoot
\bottomrule
\endlastfoot
\multicolumn{7}{l}{\textit{Panel A: direct barycentric-$\rho$ contrasts}} \\
Qwen3-4B full $-$ Qwen3-8B full & 0.00134 & [0.00058, 0.00219] & 0.001 & 0.007 & -- & -- \\
Qwen3-4B@1024 $-$ Qwen3-8B@1024 & -0.00179 & [-0.00291, -0.00077] & 0.004 & 0.008 & -- & -- \\
Qwen3-8B@1024 $-$ Qwen3-8B full & 0.00090 & [0.00001, 0.00164] & 0.049 & 0.049 & -- & -- \\
Qwen3-8B@256 $-$ Qwen3-8B@1024 & -0.00449 & [-0.00600, -0.00324] & 0.001 & 0.007 & -- & -- \\
Qwen3-8B@64 $-$ Qwen3-8B@256 & -0.01805 & [-0.02296, -0.01353] & 0.001 & 0.007 & -- & -- \\
BGE-large-v1.5 full $-$ Qwen3-8B@1024 & -0.00392 & [-0.00572, -0.00245] & 0.001 & 0.007 & -- & -- \\
BGE-large-v1.5 full $-$ Qwen3-4B@1024 & -0.00213 & [-0.00359, -0.00080] & 0.002 & 0.007 & -- & -- \\
\addlinespace[1.5pt]
ChronoBERT 2018 $-$ ChronoBERT 2017 & -0.00297 & [-0.00406, -0.00215] & 0.001 & 0.028 & 0.00572 & Yes \\
ChronoBERT 2019 $-$ ChronoBERT 2017 & -0.00240 & [-0.00367, -0.00147] & 0.001 & 0.028 & 0.00572 & Yes \\
ChronoBERT 2019 $-$ ChronoBERT 2018 & 0.00056 & [-0.00057, 0.00142] & 0.257 & 1.000 & 0.00572 & Yes \\
ChronoBERT 2020 $-$ ChronoBERT 2017 & -0.00267 & [-0.00371, -0.00171] & 0.001 & 0.028 & 0.00572 & Yes \\
ChronoBERT 2020 $-$ ChronoBERT 2018 & 0.00030 & [-0.00052, 0.00135] & 0.527 & 1.000 & 0.00572 & Yes \\
ChronoBERT 2020 $-$ ChronoBERT 2019 & -0.00026 & [-0.00134, 0.00109] & 0.693 & 1.000 & 0.00572 & Yes \\
ChronoBERT 2021 $-$ ChronoBERT 2017 & -0.00740 & [-0.00931, -0.00605] & 0.001 & 0.028 & 0.00572 & No \\
ChronoBERT 2021 $-$ ChronoBERT 2018 & -0.00444 & [-0.00577, -0.00340] & 0.001 & 0.028 & 0.00572 & No \\
ChronoBERT 2021 $-$ ChronoBERT 2019 & -0.00500 & [-0.00646, -0.00381] & 0.001 & 0.028 & 0.00572 & No \\
ChronoBERT 2021 $-$ ChronoBERT 2020 & -0.00474 & [-0.00631, -0.00352] & 0.001 & 0.028 & 0.00572 & No \\
ChronoBERT 2022 $-$ ChronoBERT 2017 & -0.00200 & [-0.00281, -0.00128] & 0.001 & 0.028 & 0.00572 & Yes \\
ChronoBERT 2022 $-$ ChronoBERT 2018 & 0.00096 & [-0.00001, 0.00214] & 0.052 & 0.364 & 0.00572 & Yes \\
ChronoBERT 2022 $-$ ChronoBERT 2019 & 0.00040 & [-0.00055, 0.00170] & 0.454 & 1.000 & 0.00572 & Yes \\
ChronoBERT 2022 $-$ ChronoBERT 2020 & 0.00066 & [-0.00044, 0.00174] & 0.242 & 1.000 & 0.00572 & Yes \\
ChronoBERT 2022 $-$ ChronoBERT 2021 & 0.00540 & [0.00408, 0.00713] & 0.001 & 0.028 & 0.00572 & No \\
ChronoBERT 2023 $-$ ChronoBERT 2017 & 0.00217 & [0.00099, 0.00367] & 0.001 & 0.028 & 0.00572 & No \\
ChronoBERT 2023 $-$ ChronoBERT 2018 & 0.00514 & [0.00359, 0.00720] & 0.001 & 0.028 & 0.00572 & No \\
ChronoBERT 2023 $-$ ChronoBERT 2019 & 0.00458 & [0.00297, 0.00673] & 0.001 & 0.028 & 0.00572 & No \\
ChronoBERT 2023 $-$ ChronoBERT 2020 & 0.00484 & [0.00345, 0.00665] & 0.001 & 0.028 & 0.00572 & No \\
ChronoBERT 2023 $-$ ChronoBERT 2021 & 0.00958 & [0.00734, 0.01257] & 0.001 & 0.028 & 0.00572 & No \\
ChronoBERT 2023 $-$ ChronoBERT 2022 & 0.00418 & [0.00286, 0.00593] & 0.001 & 0.028 & 0.00572 & No \\
ChronoBERT 2024 $-$ ChronoBERT 2017 & -0.00042 & [-0.00125, 0.00056] & 0.369 & 1.000 & 0.00572 & Yes \\
ChronoBERT 2024 $-$ ChronoBERT 2018 & 0.00255 & [0.00146, 0.00410] & 0.001 & 0.028 & 0.00572 & No \\
ChronoBERT 2024 $-$ ChronoBERT 2019 & 0.00199 & [0.00093, 0.00350] & 0.001 & 0.028 & 0.00572 & Yes \\
ChronoBERT 2024 $-$ ChronoBERT 2020 & 0.00225 & [0.00116, 0.00349] & 0.001 & 0.028 & 0.00572 & Yes \\
ChronoBERT 2024 $-$ ChronoBERT 2021 & 0.00699 & [0.00531, 0.00930] & 0.001 & 0.028 & 0.00572 & No \\
ChronoBERT 2024 $-$ ChronoBERT 2022 & 0.00159 & [0.00054, 0.00284] & 0.005 & 0.040 & 0.00572 & Yes \\
ChronoBERT 2024 $-$ ChronoBERT 2023 & -0.00259 & [-0.00361, -0.00177] & 0.001 & 0.028 & 0.00572 & Yes \\
\midrule
\multicolumn{7}{l}{\textit{Panel B: barycentric-minus-RBF gap contrasts}} \\
Qwen3-4B full $-$ Qwen3-8B full & 0.00158 & [0.00076, 0.00255] & 0.001 & 0.007 & -- & -- \\
Qwen3-4B@1024 $-$ Qwen3-8B@1024 & -0.00173 & [-0.00283, -0.00075] & 0.004 & 0.012 & -- & -- \\
Qwen3-8B@1024 $-$ Qwen3-8B full & 0.00083 & [-0.00007, 0.00159] & 0.081 & 0.162 & -- & -- \\
Qwen3-8B@256 $-$ Qwen3-8B@1024 & -0.00593 & [-0.00762, -0.00451] & 0.001 & 0.007 & -- & -- \\
Qwen3-8B@64 $-$ Qwen3-8B@256 & -0.02002 & [-0.02565, -0.01496] & 0.001 & 0.007 & -- & -- \\
BGE-large-v1.5 full $-$ Qwen3-8B@1024 & -0.00244 & [-0.00392, -0.00091] & 0.001 & 0.007 & -- & -- \\
BGE-large-v1.5 full $-$ Qwen3-4B@1024 & -0.00071 & [-0.00192, 0.00073] & 0.276 & 0.276 & -- & -- \\
\addlinespace[1.5pt]
ChronoBERT 2018 $-$ ChronoBERT 2017 & -0.00268 & [-0.00375, -0.00190] & 0.001 & 0.028 & 0.00392 & Yes \\
ChronoBERT 2019 $-$ ChronoBERT 2017 & -0.00174 & [-0.00296, -0.00081] & 0.001 & 0.028 & 0.00392 & Yes \\
ChronoBERT 2019 $-$ ChronoBERT 2018 & 0.00095 & [-0.00015, 0.00181] & 0.083 & 0.498 & 0.00392 & Yes \\
ChronoBERT 2020 $-$ ChronoBERT 2017 & -0.00203 & [-0.00301, -0.00110] & 0.001 & 0.028 & 0.00392 & Yes \\
ChronoBERT 2020 $-$ ChronoBERT 2018 & 0.00065 & [-0.00015, 0.00169] & 0.132 & 0.660 & 0.00392 & Yes \\
ChronoBERT 2020 $-$ ChronoBERT 2019 & -0.00029 & [-0.00137, 0.00105] & 0.650 & 1.000 & 0.00392 & Yes \\
ChronoBERT 2021 $-$ ChronoBERT 2017 & -0.00648 & [-0.00829, -0.00521] & 0.001 & 0.028 & 0.00392 & No \\
ChronoBERT 2021 $-$ ChronoBERT 2018 & -0.00380 & [-0.00509, -0.00280] & 0.001 & 0.028 & 0.00392 & No \\
ChronoBERT 2021 $-$ ChronoBERT 2019 & -0.00475 & [-0.00621, -0.00357] & 0.001 & 0.028 & 0.00392 & No \\
ChronoBERT 2021 $-$ ChronoBERT 2020 & -0.00445 & [-0.00603, -0.00322] & 0.001 & 0.028 & 0.00392 & No \\
ChronoBERT 2022 $-$ ChronoBERT 2017 & -0.00144 & [-0.00229, -0.00068] & 0.001 & 0.028 & 0.00392 & Yes \\
ChronoBERT 2022 $-$ ChronoBERT 2018 & 0.00124 & [0.00026, 0.00242] & 0.015 & 0.105 & 0.00392 & Yes \\
ChronoBERT 2022 $-$ ChronoBERT 2019 & 0.00030 & [-0.00066, 0.00160] & 0.613 & 1.000 & 0.00392 & Yes \\
ChronoBERT 2022 $-$ ChronoBERT 2020 & 0.00059 & [-0.00053, 0.00167] & 0.312 & 1.000 & 0.00392 & Yes \\
ChronoBERT 2022 $-$ ChronoBERT 2021 & 0.00504 & [0.00378, 0.00669] & 0.001 & 0.028 & 0.00392 & No \\
ChronoBERT 2023 $-$ ChronoBERT 2017 & 0.00293 & [0.00175, 0.00436] & 0.001 & 0.028 & 0.00392 & No \\
ChronoBERT 2023 $-$ ChronoBERT 2018 & 0.00561 & [0.00413, 0.00759] & 0.001 & 0.028 & 0.00392 & No \\
ChronoBERT 2023 $-$ ChronoBERT 2019 & 0.00466 & [0.00314, 0.00668] & 0.001 & 0.028 & 0.00392 & No \\
ChronoBERT 2023 $-$ ChronoBERT 2020 & 0.00496 & [0.00363, 0.00665] & 0.001 & 0.028 & 0.00392 & No \\
ChronoBERT 2023 $-$ ChronoBERT 2021 & 0.00941 & [0.00727, 0.01219] & 0.001 & 0.028 & 0.00392 & No \\
ChronoBERT 2023 $-$ ChronoBERT 2022 & 0.00437 & [0.00308, 0.00609] & 0.001 & 0.028 & 0.00392 & No \\
ChronoBERT 2024 $-$ ChronoBERT 2017 & 0.00026 & [-0.00058, 0.00122] & 0.548 & 1.000 & 0.00392 & Yes \\
ChronoBERT 2024 $-$ ChronoBERT 2018 & 0.00295 & [0.00189, 0.00446] & 0.001 & 0.028 & 0.00392 & No \\
ChronoBERT 2024 $-$ ChronoBERT 2019 & 0.00200 & [0.00102, 0.00342] & 0.001 & 0.028 & 0.00392 & Yes \\
ChronoBERT 2024 $-$ ChronoBERT 2020 & 0.00229 & [0.00123, 0.00346] & 0.001 & 0.028 & 0.00392 & Yes \\
ChronoBERT 2024 $-$ ChronoBERT 2021 & 0.00675 & [0.00515, 0.00894] & 0.001 & 0.028 & 0.00392 & No \\
ChronoBERT 2024 $-$ ChronoBERT 2022 & 0.00170 & [0.00062, 0.00293] & 0.004 & 0.032 & 0.00392 & Yes \\
ChronoBERT 2024 $-$ ChronoBERT 2023 & -0.00266 & [-0.00365, -0.00186] & 0.001 & 0.028 & 0.00392 & Yes \\
\end{longtable}
\endgroup


For moderate Qwen3-8B truncation from full width to 1,024 coordinates, the
direct-$\rho$ contrast narrowly rejects zero after Holm adjustment, whereas the
barycentric-minus-RBF gap interval contains zero.
The subsequent 1,024-to-256 and 256-to-64 steps yield lower fitted feedback for
the barycentric field and smaller gaps relative to RBF after Holm adjustment.
The full-width 4B cell raises both estimands, whereas the fixed-width 4B
cell lowers both after Holm adjustment; the opposing native- and fixed-width
comparisons therefore cannot be read as one pure capacity effect.
BGE-large lowers both estimands relative to Qwen3-8B@1,024, but the comparison
does not isolate architecture.

The ChronoBERT staircase is non-monotone in both estimands.
Many annual pairs reject zero after Holm adjustment, but only a subset meets
the joint equivalence rule.
The first annual step lowers both estimands after adjustment: its direct-$\rho$
and gap intervals lie within their respective equivalence bounds.
The 2017-to-2024 endpoint meets both bounds without Holm-adjusted rejection.
Thus bounded differences for selected pairs do not establish staircase-wide
invariance, an exact null, or a causal effect of later information.
Log-likelihood remains descriptive throughout.

\subsection{Annual Joint-Field Persistence}
\label{sec:p5-two-field-annual}

\Cref{tab:two-field-annual} re-estimates both coefficients within each
non-overlapping annual return sample while holding the two 2018--2022 matrices
fixed, using the raw-return, common-intercept specification of
\ppAppendixRef{sec:app-unadjusted}.

\begin{table}[htbp]
  \centering
  \caption{Annual joint estimates under the frozen barycentric and news-link
    fields in the raw-return, common-intercept specification. The 2026 period
  ends on \ppvalue{post-evaluation-date-end}.}
  \label{tab:two-field-annual}
  \begin{tabular}{lrrrr}
    \toprule
    Evaluation year & $\hat\rho^\flat$ & $\hat\rho^{\mathrm{news}}$ &
    $\hat\lambda^\flat$
    & $\hat\lambda^{\mathrm{news}}$ \\
    \midrule
    2023 & \ppnum[3]{post-eval2023-two-field-rho-b} &
    \ppnum[3]{post-eval2023-two-field-rho-n} &
    \ppnum[2]{post-eval2023-two-field-lambda-b} &
    \ppnum[2]{post-eval2023-two-field-lambda-n} \\
    2024 & \ppnum[3]{post-eval2024-two-field-rho-b} &
    \ppnum[3]{post-eval2024-two-field-rho-n} &
    \ppnum[2]{post-eval2024-two-field-lambda-b} &
    \ppnum[2]{post-eval2024-two-field-lambda-n} \\
    2025 & \ppnum[3]{post-eval2025-two-field-rho-b} &
    \ppnum[3]{post-eval2025-two-field-rho-n} &
    \ppnum[2]{post-eval2025-two-field-lambda-b} &
    \ppnum[2]{post-eval2025-two-field-lambda-n} \\
    2026 & \ppnum[3]{post-eval2026-two-field-rho-b} &
    \ppnum[3]{post-eval2026-two-field-rho-n} &
    \ppnum[2]{post-eval2026-two-field-lambda-b} &
    \ppnum[2]{post-eval2026-two-field-lambda-n} \\
    \bottomrule
  \end{tabular}
\end{table}

Both channel coefficients remain positive in every year, and the barycentric
estimate is larger throughout.
The gap narrows in the incomplete 2026 period, which uses only
\ppvalue{post-eval2026-two-field-trading-days} dates and is not directly
comparable with a complete-year estimate.

\section{Formal Verification Status}\label{sec:lean-status}

Selected algebraic closure and scalar-projection results are machine-checked
in Lean 4 with mathlib under the hypotheses recorded below.
The factor-beta invariance proposition is not among these
machine-checked statements.
These checks concern the mathematical statements; they do not certify the
text-derived interaction weights or the statistical interpretation of the
estimates, which depends on the empirical design in \Cref{sec:empirical-design}.
The scripts and theorem inventory are included in the replication archive.

\ppLeanStatusList
